% ---- FIGURE: circuit -> architecture graph -> lattice surgery ----
\begin{figure*}[t]
    \centering
    \resizebox{0.8\textwidth}{!}{%
    \begin{tikzpicture}
      % ======================== (a) circuit ========================
      \begin{scope}
        \foreach \y/\name in {1.4/q_0,0.7/q_1,0/q_2}{
            \draw[ftgray] (0,\y)--(2,\y);
            \node[left,qlabel] at (-0.05,\y){$\name$};}
        % CNOT(q0,q1) at x=0.6
        \fill (0.6,1.4) circle (0.05);
        \draw (0.6,1.4)--(0.6,0.7);
        \draw (0.6,0.7) circle (0.12);
        \draw (0.48,0.7)--(0.72,0.7) (0.6,0.58)--(0.6,0.82);
        % T on q2 at x=0.6
        \node[draw,fill=white,inner sep=1.2pt,font=\footnotesize] at (0.6,0){$T$};
        % CNOT(q1,q2) at x=1.4
        \fill (1.4,0.7) circle (0.05);
        \draw (1.4,0.7)--(1.4,0);
        \draw (1.4,0) circle (0.12);
        \draw (1.28,0)--(1.52,0) (1.4,-0.12)--(1.4,0.12);
        \node[ftgray,font=\footnotesize] at (1,-0.7){(a) circuit $\C$};
      \end{scope}

      % map arrow
      \draw[-{Stealth[length=2mm]},thick,ftgray]
            (2.25,0.7)--(3.15,0.7) node[midway,above,font=\footnotesize,black]{map};

      % ==================== (b) architecture graph ====================
      \begin{scope}[xshift=3.6cm]
        \foreach \y in {0,1,2}\draw[ftgray!40] (0,\y)--(2,\y);
        \foreach \x in {0,1,2}\draw[ftgray!40] (\x,0)--(\x,2);
        \foreach \x in {0,1,2}\foreach \y in {0,1,2}\node[free] at (\x,\y){};
        \node[data] at (0,2){}; \node[qlabel] at (0,2){$q_0$};
        \node[data] at (0,0){}; \node[qlabel] at (0,0){$q_1$};
        \node[data] at (2,0){}; \node[qlabel] at (2,0){$q_2$};
        \node[magic] at (2,2){}; \node[mstar] at (2,2){};
        \node[ftgray,font=\footnotesize] at (1,-0.7){(b) mapping on $\A$};
      \end{scope}

      % route arrow
      \draw[-{Stealth[length=2mm]},thick,ftgray]
            (6.1,0.7)--(7.1,0.7) node[midway,above,font=\footnotesize,black]{route};

      % ==================== (c) lattice-surgery step ====================
      \begin{scope}[xshift=7.5cm]
        \foreach \y in {0,1,2}\draw[ftgray!40] (0,\y)--(2,\y);
        \foreach \x in {0,1,2}\draw[ftgray!40] (\x,0)--(\x,2);
        \foreach \x in {0,1,2}\foreach \y in {0,1,2}\node[free] at (\x,\y){};
        % corridors of the first layer: CNOT(q0,q1) and T(q2)
        \node[corridor] at (0,1){};                        % CNOT(q0,q1)
        \node[tile,fill=ftorange!22,draw=ftorange!55!black] at (2,1){}; % T route
        % patches
        \node[data] at (0,2){}; \node[qlabel] at (0,2){$q_0$};
        \node[data] at (0,0){}; \node[qlabel] at (0,0){$q_1$};
        \node[data] at (2,0){}; \node[qlabel] at (2,0){$q_2$};
        \node[magic] at (2,2){}; \node[mstar] at (2,2){};
        % merge (CNOT q0-q1) and T-route (q2 -> magic) highlights
        \draw[ftgreen!55!black,line width=1.1pt,shorten >=2.5mm,shorten <=2.5mm]
              (0,2)--(0,0);
        \draw[pin,ftorange!70!black] (2,0)--(2,2);
        \node[ftgray,font=\footnotesize] at (1,-0.7){(c) first layer};
      \end{scope}
    \end{tikzpicture}}
    \caption{From circuit to lattice surgery. A Clifford$+T$ circuit~(a) is
    \emph{mapped} onto the architecture graph $\A$~(b), assigning each logical
    qubit to a tile while magic states ($\star$) stay fixed. The \emph{router}
    then executes each layer by reserving vertex-disjoint corridors; panel~(c)
    shows the first layer, where a green merge corridor realises $\CNOT(q_0,q_1)$
    and an orange lane routes the $T$ gate on $q_2$ to a magic state. Choosing the
    most gates servable at once is a maximum vertex-disjoint paths problem on
    $\A$.}
    \label{fig:overview}
\end{figure*}
